% ============================================================
% style/headings.tex
%
% Оформление стандартных заголовков документа.
%
% Дополнительная цель файла:
%   не засорять дерево структуры LaTeX Workshop служебными строками
%   из \titleformat и \titlespacing.
%
% Поэтому стандартные команды заголовков внутри настроек вызываются
% через \csname ...\endcsname, а не напрямую.
% ============================================================


% ------------------------------------------------------------
% 1. Глубина нумерации и оглавления
% ------------------------------------------------------------
%
% secnumdepth=1:
%   - главы нумеруются;
%   - разделы нумеруются;
%   - подразделы не нумеруются.
%
% tocdepth=1:
%   - в оглавление попадают главы и разделы;
%   - подразделы в оглавление не попадают.

\setcounter{secnumdepth}{1}
\setcounter{tocdepth}{1}


% ------------------------------------------------------------
% 2. Внутренние имена для numberless-форматов titlesec
% ------------------------------------------------------------
%
% titlesec для заголовков без номера использует синтаксис:
%
%   name=<команда>,numberless
%
% Чтобы LaTeX Workshop не принимал служебные настройки за настоящие
% заголовки документа, имена команд собираются через \csname.

\expandafter\edef\csname rPartName\endcsname{%
    \expandafter\string\csname part\endcsname%
}

\expandafter\edef\csname rChapterName\endcsname{%
    \expandafter\string\csname chapter\endcsname%
}

\expandafter\edef\csname rSectionName\endcsname{%
    \expandafter\string\csname section\endcsname%
}

\expandafter\edef\csname rSubsectionName\endcsname{%
    \expandafter\string\csname subsection\endcsname%
}

\expandafter\edef\csname rSubsubsectionName\endcsname{%
    \expandafter\string\csname subsubsection\endcsname%
}


% ------------------------------------------------------------
% 3. Нумерация задач внутри главы
% ------------------------------------------------------------
%
% По умолчанию в report было бы:
%
%   1.1, 1.2, 1.3
%
% Здесь будет:
%
%   1, 2, 3
%
% При новой главе счётчик разделов сбрасывается автоматически.

\expandafter\renewcommand\csname thesection\endcsname{%
    \arabic{section}%
}


% ------------------------------------------------------------
% 4. Оформление частей
% ------------------------------------------------------------

\expandafter\titleformat\expandafter{\csname part\endcsname}[display]
    {\thispagestyle{empty}\normalfont\centering\color{rtext}}
    {
        {\Large\bfseries Часть \thepart}
    }
    {1.2em}
    {
        \Huge\bfseries
    }
    []

\expandafter\titlespacing\expandafter*\expandafter{\csname part\endcsname}
    {0pt}
    {0.26\textheight}
    {0pt}


% ------------------------------------------------------------
% 5. Оформление глав
% ------------------------------------------------------------
%
% Вид:
%
%                                   Глава 1
%
%                              Название главы
%
% Номер главы справа, название ниже.
% Линия между номером и названием не используется.

\expandafter\titleformat\expandafter{\csname chapter\endcsname}[display]
    {\normalfont\filleft\color{rtext}}
    {
        {\Large\bfseries Глава \thechapter}
    }
    {0.5em}
    {
        \Huge\bfseries
    }
    []

\expandafter\titlespacing\expandafter*\expandafter{\csname chapter\endcsname}
    {0pt}
    {-4mm}
    {14mm}


% ------------------------------------------------------------
% 6. Оформление глав без номера
% ------------------------------------------------------------

\titleformat{name=\rChapterName,numberless}[display]
    {\normalfont\filleft\color{rtext}}
    {}
    {0pt}
    {
        \Huge\bfseries
    }
    []

\titlespacing*{name=\rChapterName,numberless}
    {0pt}
    {-4mm}
    {14mm}


% ------------------------------------------------------------
% 7. Подпись под главой
% ------------------------------------------------------------
%
% Использование в документе:
%
%   \rsubtitle{Первый контакт}
%
% Это не структурная команда, а декоративная подпись.

\NewDocumentCommand{\rsubtitle}{m}{%
    \vspace{-10mm}%
    \begin{flushright}
        {\large\itshape\color{rtext}#1}%
    \end{flushright}%
    \vspace{8mm}%
}


% ------------------------------------------------------------
% 8. Оформление разделов
% ------------------------------------------------------------
%
% Использование в документе:
%
%   раздел = задача
%
% Автонумерация:
%
%   1 Задача
%   2 Задача
%   3 Задача

\expandafter\titleformat\expandafter{\csname section\endcsname}[hang]
    {\normalfont\Large\bfseries\color{rtext}}
    {
        \csname thesection\endcsname
    }
    {0.8em}
    {}
    []

\expandafter\titlespacing\expandafter*\expandafter{\csname section\endcsname}
    {0pt}
    {1.8\baselineskip}
    {0.75\baselineskip}


% ------------------------------------------------------------
% 9. Оформление разделов без номера
% ------------------------------------------------------------

\titleformat{name=\rSectionName,numberless}[hang]
    {\normalfont\Large\bfseries\color{rtext}}
    {}
    {0pt}
    {}
    []

\titlespacing*{name=\rSectionName,numberless}
    {0pt}
    {1.8\baselineskip}
    {0.75\baselineskip}


% ------------------------------------------------------------
% 10. Оформление подразделов
% ------------------------------------------------------------
%
% Формат runin означает, что заголовок идёт в одну строку
% с последующим текстом.
%
% Пример визуального вида:
%
%   Условие -- здесь начинается текст условия.

\expandafter\titleformat\expandafter{\csname subsection\endcsname}[block]
    {\normalfont\bfseries\color{rtext}}
    {}
    {0pt}
    {}

\expandafter\titlespacing\expandafter*\expandafter{\csname subsection\endcsname}
    {0pt}
    {0.8\baselineskip}
    {0.35\baselineskip}


% ------------------------------------------------------------
% 11. Оформление подразделов без номера
% ------------------------------------------------------------

\titleformat{name=\rSubsectionName,numberless}[block]
    {\normalfont\bfseries\color{rtext}}
    {}
    {0pt}
    {}

\titlespacing*{name=\rSubsectionName,numberless}
    {0pt}
    {0.8\baselineskip}
    {0.35\baselineskip}

% ------------------------------------------------------------
% 12. Оформление малых внутренних заголовков
% ------------------------------------------------------------
%
% Они тоже идут в строку с последующим текстом, но набираются курсивом.

\expandafter\titleformat\expandafter{\csname subsubsection\endcsname}[block]
    {\normalfont\itshape\color{rtext}}
    {}
    {0pt}
    {}
    []

\expandafter\titlespacing\expandafter*\expandafter{\csname subsubsection\endcsname}
    {0pt}
    {0.6\baselineskip}
    {0.45em}


% ------------------------------------------------------------
% 13. Оформление малых внутренних заголовков без номера
% ------------------------------------------------------------

\titleformat{name=\rSubsubsectionName,numberless}[block]
    {\normalfont\itshape\color{rtext}}
    {}
    {0pt}
    {}
    [  ]

\titlespacing*{name=\rSubsubsectionName,numberless}
    {0pt}
    {0.6\baselineskip}
    {0.45em}